\chapter{Relative Donaldson--Thomas operators}
\label{ch:relative-dt}

\section{Boundary amplitudes and the Euler index}
\label{rdt:sec:index}

A closed curve-counting series forgets the dependence on boundary
insertions. Relative ideal sheaves retain those insertions: restricting
an ideal to a relative K3 fibre gives a point of a Hilbert scheme. The
Poincare diagonal then provides a geometric way to join two ends.
The resulting operator is determined before its trace is taken.

Let $S$ be a smooth projective complex K3 surface, let $E$ be a smooth
elliptic curve, and let $0\ne\beta\in\operatorname{Pic}(S)$ be effective.
The integer $d\ge0$ denotes degree over the base curve. Put
\[
 Y_d=S^{[d]},\qquad H_d=H^*(Y_d,\Q),\qquad
 \eta_d(a,b)=\int_{Y_d}a\cup b.
\]
We use the actual cohomological grading, without a shift. Thus $Y_d$ has
complex dimension $2d$, the pairing has degree $-4d$, and $H_d$ is
concentrated in even degrees. The last assertion follows, for example,
from the Nakajima description in terms of the even cohomology of $S$.
For $d=0$, $Y_0$ is a point and $H_0=\Q$.

\begin{convention}\label{rdt:conv:index}
The notation $(\beta,n,m,p)$ in \cite{OberdieckMC} becomes
$(\beta,d,m,r)$ here. On $S\times E$, its $m$ is the Euler
characteristic denoted by $n$ in the closed ideal-sheaf charge.
On a relative geometry over a general curve, however, the paper defines
$m=\int\ch_3(\mathcal O_Z)$, not the Euler characteristic.
Write $N=\chi(\mathcal O_Z)$ for the Euler characteristic on the tube.
\end{convention}

\begin{lemma}[The sewing index]\label{rdt:lem:index}
For a curve of class $(\beta,d)$ on $S\times A$, with $A$ smooth of
genus $g$, one has
\begin{equation}\label{rdt:eq:RR}
 \chi(\mathcal O_Z)=m+d(1-g).
\end{equation}
Consequently $N=m+d$ on $S\times\mathbb P^1$, whereas the closed
Euler characteristic on $S\times E$ is $n=m$.
\end{lemma}
\begin{proof}
The rank and first Chern character of $\mathcal O_Z$ vanish. Its
Riemann--Roch formula is therefore
\[
 \chi(\mathcal O_Z)=\int\ch_3(\mathcal O_Z)
   +\frac12\int[Z]c_1(T_{S\times A}).
\]
The K3 tangent bundle has first Chern class zero. The horizontal degree
contributes $d\deg c_1(T_A)/2=d(1-g)$, proving the formula. On expanded
relative targets the conserved numerical index is defined by the same
contraction convention; the normalization exact sequence at a node
gives the identical relation.
\end{proof}

If a nodal curve is separated into two components, admissibility of the
relative ideal gives
\begin{equation}\label{rdt:eq:normalization}
 0\longrightarrow\mathcal O_Z\longrightarrow
 \mathcal O_{Z_1}\oplus\mathcal O_{Z_2}
 \longrightarrow\mathcal O_\eta\longrightarrow0,
 \qquad \operatorname{length}(\eta)=d.
\end{equation}
Thus $N=N_1+N_2-d$. In terms of $m$, the same identity is $m=m_1+m_2$.
For a nonseparating node it is $n=N-d$. This is the reason for the
indexing in \cite[Remark 5.1]{OberdieckMC}.

\subsection{Actual charges and saturation}

On the closed Calabi--Yau product, the ideal sequence and
$\sqrt{\operatorname{td}(S)}=1+[\mathrm{pt}]$ give
\begin{equation}\label{rdt:eq:charges}
 P=(1,0,1-d),\qquad Q=(0,-\beta,-n),\qquad
 \Pi_X(Q,P)=(\beta^2/2,n,d-1).
\end{equation}
On a tube the same surface-valued Chern character has $m=N-d$ in the
last component of $Q$. It would be incorrect to put the unshifted
tube Euler characteristic in that component.

\begin{proposition}\label{rdt:prop:wedge}
For the pair \eqref{rdt:eq:charges}, its rank-two saturation index is
\[
 [\operatorname{Sat}(\Z P+\Z Q):\Z P+\Z Q]=\gcd(\operatorname{div}(\beta),n).
\]
The gcd of the entries of its Gram matrix is instead
$\gcd(\beta^2,2(d-1),n)$.
\end{proposition}
\begin{proof}
Write $\beta=\sum b_i e_i$ in an integral basis. The two-column
matrix with columns $P,Q$ has minors containing every $b_i$ and $n$;
the remaining possibly nonzero minors are $(1-d)b_i$ up to sign.
The second determinantal divisor is therefore $\gcd(b_i,n)$, which is
the saturation index by Smith normal form. The first coordinate of
$P$ is one, so the rank is two and no additional first determinantal
divisor occurs. The Mukai pairings are $Q^2=\beta^2$,
$P^2=2(d-1)$ and $Q\cdot P=n$, proving the other formula.
\end{proof}

For a section $C$ and fibre $F$ of an elliptic K3,
$C^2=-2$, $F^2=0$, $C\cdot F=1$. In particular $C$ is primitive.
For $d=1$, the following table uses the \emph{closed} Euler index.
\begin{center}
\begin{tabular}{ccccc}
\toprule
$\beta$ & $n$ & $\operatorname{div}(\beta)$ & wedge index & Gram gcd\\
\midrule
$C$ & 1 & 1 & 1 & 1\\
$C$ & 2 & 1 & 1 & 2\\
$2C$ & 1 & 2 & 1 & 1\\
$2C$ & 2 & 2 & 2 & 2\\
$2C$ & 0 & 2 & 2 & 8\\
\bottomrule
\end{tabular}
\end{center}
Neither of these arithmetic invariants is a signed BKM root
multiplicity. Nor does a saturated formal pair with the same Gram
matrix give a replacement for the actual ideal-sheaf charge.

\section{The evaluation correspondence}
\label{rdt:sec:kernel}

Let $M_{\beta,d,m}(A;\mathbf z)$ be the stable expanded relative
ideal-sheaf moduli for
\[
 S\times A\big/\bigcup_{z\in\mathbf z}S_z
\]
with class $(\beta,d)$ and third Chern character $m$.
Admissibility means normality to the distinguished divisors and to the
nodes of the expansion; stability requires finite automorphisms of the
expanded object. The relative compactification is proper for this
projective target. Each restriction is a length-$d$ ideal on $S$, giving
proper evaluation maps $\operatorname{ev}_z:M\to Y_d$.
We use the relative obstruction theory and reduced degeneration theorem
in the form of \cite[\S5.1]{OberdieckMC}, based on
\cite{LiWu,MaulikPandharipandeThomas}. These virtual-theoretic inputs are
part of the definition of the relative DT invariants used below.

Rational Chow groups retain the finite automorphism weights of the
expanded moduli. The ordinary virtual dimension is $(2-2g(A))d$.
For $(\beta,m)\ne(0,0)$, removal of the K3 cosection gives the
reduced virtual class of dimension $(2-2g(A))d+1$. We write
$[M]^\mathrm{red}$ for the class denoted $[M]^{\mathrm{vir}}$ in the paper,
and $[M]^\mathrm{std}$ for its ordinary class.

On the universal expanded target let $\pi$ be projection to the moduli,
$c$ the contraction to $S\times A$, and $\mathcal Z$ the universal
subscheme. The descendant convention is
\begin{equation}\label{rdt:eq:desc}
 \tau_i(\gamma)=\pi_*\bigl(\ch_{2+i}(\mathcal O_{\mathcal Z})c^*\gamma\bigr).
\end{equation}
For an interior point class $\omega$ and a divisor $D$ on $S$,
$\tau_0(\omega D)$ has real degree two. The convention uses
$\mathcal O_{\mathcal Z}$: in terms of its ideal sheaf it is minus
the corresponding $\ch_2(\mathcal I)$ expression. There is no additional
sign attached to the descendant.

\begin{definition}[Marked tube kernel]\label{rdt:def:kernel}
For $A=\mathbb P^1$ and $\mathbf z=(0,\infty)$, define
\begin{equation}\label{rdt:eq:kernel}
 K^D_{\beta,d,m}=(\operatorname{ev}_0,\operatorname{ev}_\infty)_*
 \bigl(\tau_0(\omega D)\cap[M_{\beta,d,m}]^\mathrm{red}\bigr)
 \in A_{2d}(Y_d\times Y_d)_\Q.
\end{equation}
Its homology degree is $4d$. Its Poincare dual $\kappa^D_{\beta,d,m}$
has cohomology degree $4d$. The associated degree-zero endomorphism is
\begin{equation}\label{rdt:eq:B}
 \widehat{\mathsf B}^D_{\beta,d,m}(a)
 =\operatorname{pr}_{2*}
   \bigl(\operatorname{pr}_1^*a\cup\kappa^D_{\beta,d,m}\bigr).
\end{equation}
Equivalently,
\begin{equation}\label{rdt:eq:matrix}
 \eta_d(\widehat{\mathsf B}^D_{\beta,d,m}a,b)
 =\int_{[M_{\beta,d,m}]^\mathrm{red}}
    \tau_0(\omega D)\operatorname{ev}_0^*a\operatorname{ev}_\infty^*b.
\end{equation}
\end{definition}

The virtual dimension is $2d+1$, so the descendant in
\eqref{rdt:eq:kernel} leaves dimension $2d$, as required for a
degree-zero correspondence on $Y_d$. Properness supplies every
pushforward. Nondegeneracy of $\eta_d$ proves existence and uniqueness
of the endomorphism from its matrix coefficients. Formula
\eqref{rdt:eq:kernel}, rather than a choice of a matrix with prescribed
trace, is its geometric definition.

For $D$ algebraic, the kernel is a rational algebraic cycle class.
The same definition with a real class $D\in H^2(S,\R)$ uses the
cohomological descendant and gives a real homology class and
endomorphism. It need not give an algebraic Chow class in that case.

\begin{proposition}[Laurent completion]\label{rdt:prop:Laurent}
The series
\begin{equation}\label{rdt:eq:series}
 \widehat{\mathsf B}^D_{\beta,d}(r)=\sum_{m\in\Z}(-r)^m\widehat{\mathsf B}^D_{\beta,d,m}
\end{equation}
belongs to $\End_\Q(H_d)((r))$. A lower bound for nonzero matrix
coefficients is
\begin{equation}\label{rdt:eq:lower}
 m\ge L_{\beta,d}:=\min_{k\mid\beta}
       \left(-kd-\frac{\beta^2}{2k}\right).
\end{equation}
\end{proposition}
\begin{proof}
For a cap, the localization proof in \cite[\S\S4.3--5.3]{OberdieckMC}
eliminates mixed fixed components and all extremal flags with more
than one nontrivial step. Pure-rubber cap terms vanish as well:
the non-equivariant insertion degree is $2d+1$, whereas their
reduced virtual dimension is $2d$. A remaining contribution is indexed by
$k\mid\beta$, $k\mid m$, and a nested Hilbert scheme with length
\[
 d_1=d+\frac{m}{k}+\frac{\beta^2}{2k^2}.
\]
A negative length is empty. Therefore that term requires
$m\ge-kd-\beta^2/(2k)$. There are finitely many divisors of the
nonzero integral class $\beta$.

Section 5.2 of the same paper expresses each fixed relative insertion
invariant as a finite $\Q$-linear combination of cap invariants.
Nonseparating degeneration lowers genus; separating a relative marking
uses an invertible ordinary descendant--Nakajima matrix. All ordinary
factors have $(\beta,m)=(0,0)$ and are independent of $r$. The finite
linear induction therefore preserves the displayed lower bound.
Applying this to a finite cohomology basis proves the assertion for
the whole operator, not merely for its trace.
\end{proof}

One may sum over effective classes in the Novikov completion defined
by an ample degree, allowing the Laurent bound in $r$ to depend on
$\beta$. At a fixed class all matrices are finite dimensional. There
are finitely many effective classes of bounded ample degree on a fixed
projective K3: an irreducible curve has square at least $-2$, and the
Hodge index theorem bounds its negative-definite component when its
ample degree is bounded; an effective divisor has a bounded number of
such components. This gives the usual locally finite Novikov
multiplication. No rationality of every matrix entry is required for
\eqref{rdt:eq:series}.

For the tube indexed literally by its Euler characteristic, put
\begin{equation}\label{rdt:eq:Eu}
 \mathsf B^{D,\mathrm{Eu}}_{\beta,d,N}:=\widehat{\mathsf B}^D_{\beta,d,N-d},
 \qquad
 \mathsf B^{D,\mathrm{Eu}}_{\beta,d}(r)=(-r)^d\widehat{\mathsf B}^D_{\beta,d}(r).
\end{equation}
This conversion will be retained in both the sewing and
multiple-cover formulas.

\section{Cutting, gluing and the elliptic quotient}
\label{rdt:sec:glue}

At a separating node, the two relative evaluations form a fibre
product $M_1\times_{Y_d}M_2$. Admissible ideals with the same boundary
ideal $I_\eta$ glue by the kernel of their difference map to
$I_\eta$. Cutting is restriction to the normalization; the absence of
the unwanted Tor term is exactly the relative admissibility condition.
This constructs the gluing morphism from the fibre product to the
central relative moduli. The virtual intersection on the fibre product
is the refined pullback by the diagonal $\Delta_{Y_d}$.

Write $\mathcal A^{\mathrm{red}}_{A}(I\mid\lambda)$ for the relative
invariant with stationary descendants $I$, relative insertions
$\lambda$, and generating variable $(-r)^m$, and similarly for the
ordinary class. With $\eta^{ij}$ the inverse pairing matrix, the
reduced degeneration formula needed here is the insertion identity
\begin{equation}\label{rdt:eq:virtualglue}
\begin{aligned}
 \mathcal A^{\mathrm{red}}_{A}(I\mid\lambda)
  =\sum_{i,j}\eta^{ij}\bigl(&
    \mathcal A^{\mathrm{red}}_{A_1}(I_1\mid\lambda_1,e_i)
    \mathcal A^{\mathrm{std}}_{A_2}(I_2\mid\lambda_2,e_j)\\
   &+\mathcal A^{\mathrm{std}}_{A_1}(I_1\mid\lambda_1,e_i)
    \mathcal A^{\mathrm{red}}_{A_2}(I_2\mid\lambda_2,e_j)\bigr).
\end{aligned}
\end{equation}
The numerical splittings are understood, and the relative and
interior marks are distributed as specified. The refined diagonal
pullback above explains the contraction in this identity. Equality
after arbitrary relative insertions is sufficient for all
cohomological correspondences below; no stronger identification of
uninserted Chow classes is needed. The standard class vanishes unless
$(\beta,m)=(0,0)$, by the K3 cosection. Thus for fixed nonzero $\beta$
exactly one factor carries the reduced class and the nonzero degree.
There is no term with two reduced factors. This is the precise
formula stated in \cite[\S5.1]{OberdieckMC}.

If $\{e_i\}$ is a homogeneous basis and $\{e^i\}$ its Poincare dual,
the coevaluation and evaluation maps are
\[
 \operatorname{coev}(1)=\sum_i e_i\otimes e^i=[\Delta_{Y_d}],
 \qquad \operatorname{ev}(a\otimes b)=\eta_d(a,b).
\]
In an arbitrary basis, coevaluation is
$\sum_{i,j}\eta^{ij}e_i\otimes e_j$.
These are genuine boundary maps on $H_d$. Partition multiplicities in
a Nakajima basis are already contained in $\eta^{ij}$.

\begin{theorem}[Geometric closure]\label{rdt:thm:trace}
Assume $\beta\cdot D\ne0$. With the quotient-stack reduced invariant
of $S\times E$,
\begin{equation}\label{rdt:eq:trace}
 \frac{1}{\beta\cdot D}\operatorname{Tr}_{H_d}\widehat{\mathsf B}^D_{\beta,d}(r)
 =\sum_{n\in\Z}(-r)^n\mathrm{DT}^{S\times E}_{n,(\beta,d)}.
\end{equation}
For a literal Euler-indexed tube this becomes
\begin{equation}\label{rdt:eq:EuTrace}
 \frac{(-r)^{-d}}{\beta\cdot D}
 \operatorname{Tr}_{H_d}\mathsf B^{D,\mathrm{Eu}}_{\beta,d}(r)
 =\sum_n(-r)^n\mathrm{DT}^{S\times E}_{n,(\beta,d)}.
\end{equation}
\end{theorem}
\begin{proof}
Degenerate $E$ to a one-nodal rational curve with the interior mark
away from the node. The normalization has the two relative fibres
$0,\infty$. The nonseparating version of
\eqref{rdt:eq:virtualglue} inserts the full diagonal in the two
evaluations. By \eqref{rdt:eq:matrix} its contraction is
\[
 \sum_i\eta_d(\widehat{\mathsf B}^D_{\beta,d}(r)e_i,e^i)
 =\operatorname{Tr}\widehat{\mathsf B}^D_{\beta,d}(r).
\]
This computes the \emph{unquotiented} reduced invariant with
$\tau_0(\omega D)$. All degrees of $H_d$ are even, so the categorical
supertrace is the ordinary trace. The fixed degeneration has labeled
normalization branches. The formula has no additional factor $1/2$;
it is not integration over the boundary of a moduli space of unlabeled
curves. Expanded-target automorphisms are already counted by the
rational virtual pushforward.

It remains to divide out elliptic translation. For $\beta\ne0$ the
translation stabilizers in the Hilbert moduli are finite. Indeed, a
positive-dimensional subgroup of an elliptic curve is the whole curve;
an invariant one-cycle would then consist of elliptic fibres over
points of $S$ and would have K3 class zero. The quotient stack can
nevertheless have nontrivial finite stabilizers.

The reduced virtual descent identity of
\cite[\S6.2]{OberdieckMC}, from \cite{OberdieckReducedPT}, is
\[
 Z^{S\times E}(\tau_0(\omega D))
   =(\beta\cdot D)\,Z^{S\times E/E}.
\]
Its normalization is also visible on an orbit. Pulling the descendant
to the translation parameter $E$ gives degree $\beta\cdot D$.
For an orbit with stabilizer $G$, the degree on $E/G$ is
$(\beta\cdot D)/|G|$, matching the quotient-stack mass $1/|G|$.
Virtual descent and the projection formula make this orbit computation
an equality of invariants. This uses the quotient virtual theory,
rather than an assertion that the action on the moduli is free.
Combining the two identities proves \eqref{rdt:eq:trace}.
Equation \eqref{rdt:eq:EuTrace} follows from
\eqref{rdt:eq:Eu}, or directly from $n=N-d$ in
\eqref{rdt:eq:normalization}.
\end{proof}

The empty relative boundary for $d=0$ is included in this proof.
It is the state $\Q$ attached to $Y_0=\mathrm{pt}$. It has no
automatic identification with a relative Hall unit $BE$ or a prepared
protected vacuum.

\section{The ordinary and reduced sewing system}
\label{rdt:sec:firstorder}

In genus zero and degree $(\beta,m)=(0,0)$, the relative moduli is
$Y_d$, with ordinary virtual class $[Y_d]$ and identity evaluation
maps \cite[Lemma 5.2]{OberdieckMC}. Consequently the ordinary cap,
tube and three-boundary operation are respectively
\[
 1\in H_d,\qquad \id_{H_d},\qquad (a,b)\longmapsto a\cup b.
\]
The ordinary tube has actual Euler characteristic $d$, although its
$m$-series is constant. The ordinary cap with $\tau_0(\omega D)$
gives the class
\[
 D^{[d]}=\tau_0^{Y_d}(D)=\pi_*\bigl(\ch_2(\mathcal O_{\mathcal U})
                      \pi_S^*D\bigr)\in H^2(Y_d),
\]
where $\mathcal U\subset Y_d\times S$ is the universal subscheme.
Its marked tube is cup multiplication $M_{D^{[d]}}$.

\begin{proposition}[The first-order Frobenius structure]
\label{rdt:prop:firstorder}
The unmarked reduced cap and tube vanish. Let $\mu_{\beta,d}(r)$
be defined by the unmarked three-relative invariant:
\[
 \eta_d(\mu_{\beta,d}(a,b),c)
   =Z^{\mathrm{red},S\times\mathbb P^1/S_0,S_1,S_\infty}_{\beta,d}(a,b,c).
\]
It has degree $-2$, is commutative and cyclic, and satisfies
\begin{align}
 \mu_{\beta,d}(1,a)&=0,\label{rdt:eq:redunit}\\
 \mu_{\beta,d}(ab,c)+\mu_{\beta,d}(a,b)c
  &=\mu_{\beta,d}(a,bc)+a\mu_{\beta,d}(b,c).
 \label{rdt:eq:assoc}
\end{align}
Over the coefficientwise Laurent/Novikov ring $\Lambda$, the product
\begin{equation}\label{rdt:eq:star}
 a\star b=ab+\epsilon\sum_{\beta>0}Q^\beta\mu_{\beta,d}(a,b),
 \qquad \epsilon^2=0,\quad |\epsilon|=2,
\end{equation}
is a unital first-order Frobenius product with unchanged Poincare
pairing and cup unit.
\end{proposition}
\begin{proof}
The reduced cap has virtual dimension $2d+1$ and its pushforward to
$Y_d$ is zero by dimension. If $R_{\beta,d}$ denotes the unmarked
reduced tube, a separating degeneration of that tube gives
\[
 R_{\beta,d}=R_{\beta,d}\id+\id R_{\beta,d}=2R_{\beta,d}.
\]
Therefore $R_{\beta,d}=0$ over $\Q$. Attaching the ordinary cap to
the three-relative operation and using the vanishing reduced cap
identifies $\mu(1,a)$ with $R(a)$, proving
\eqref{rdt:eq:redunit}.

The virtual dimension of the three-relative space is still $2d+1$,
so its trilinear pairing is nonzero only when the real degrees sum to
$4d+2$. This gives degree $-2$ for $\mu$. Permuting the relative
points gives cyclicity and symmetry. Applying reduced degeneration
to a four-relative sphere in its two channels gives exactly
\eqref{rdt:eq:assoc}: each channel has one ordinary and one reduced
factor. Expansion modulo $\epsilon^2$ proves associativity of
\eqref{rdt:eq:star}; the other stated properties follow from cyclicity
and \eqref{rdt:eq:redunit}. Its total degree is zero because
$|\epsilon|=2$.
\end{proof}

\begin{proposition}[The marked state]\label{rdt:prop:markedstate}
Let $c^D_{\beta,d}(r)\,1$ denote the pushforward of the reduced cap
with $\tau_0(\omega D)$. It is a scalar multiple of $1$ by dimension.
Then
\begin{equation}\label{rdt:eq:markeddecomposition}
 \widehat{\mathsf B}^D_{\beta,d}(a)
   =\mu_{\beta,d}(D^{[d]},a)+c^D_{\beta,d}a.
\end{equation}
In particular $\widehat{\mathsf B}^D$ is self-adjoint for $\eta_d$ and
$\widehat{\mathsf B}^D(1)=c^D\,1$. The degree-two first-order marked operator
\[
 M_{D^{[d]}}+\epsilon\sum_{\beta>0}Q^\beta\widehat{\mathsf B}^D_{\beta,d}
\]
is multiplication for $\star$ by
$D^{[d]}+\epsilon\sum_{\beta>0}Q^\beta c^D_{\beta,d}\,1$.
\end{proposition}
\begin{proof}
Degenerate the marked tube by moving the interior mark to a third cap.
When the three-relative component is reduced, the cap is ordinary and
inserts $D^{[d]}$. When the cap is reduced, the three-relative
component is ordinary and multiplies by $c^D1$. These are precisely
the two terms in \eqref{rdt:eq:markeddecomposition}. Cyclicity and the
unit identity from Proposition~\ref{rdt:prop:firstorder} prove the
remaining statements.
\end{proof}

All sewings follow the same tensor contraction rule. For example the
handle element is $\sum_i e_i\star e^i$; its ordinary part is
$\sum_i e_i e^i=c_{2d}(T_{Y_d})$. Closing an ordinary tube gives
$\operatorname{Tr}\id=\chi_{\mathrm{top}}(Y_d)$, and closing the marked reduced
coefficient gives Theorem~\ref{rdt:thm:trace}. This is a system of
curve-degeneration operations, not an identification with Hall
convolution or chiral collisions.

The reduced coefficient does not inherit a normalized-unit value.
Its unmarked unit correction is zero; its marked value is the computed
$c^D$. Likewise, for a unital map $I=I_0+\epsilon I_1$ to dual
numbers, multiplicativity gives
$I_1(ab)=I_0(a)I_1(b)+I_1(a)I_0(b)$ and $I_1(1)=0$.
The similarity to the reduced Hall integration of
\cite{OberdieckShenReduced} is an algebraic pattern. Identifying the two
systems requires a map of their actual carriers and operations.

\section{Multiple covers before taking the trace}
\label{rdt:sec:covers}

For every positive $k\mid\beta$, choose a K3 surface $S_k$ and a real
isometry
\[
 \phi_k:H^2(S,\R)\longrightarrow H^2(S_k,\R),\qquad
 \beta_k=\phi_k(\beta/k)
\]
with $\beta_k$ primitive, integral and effective, as in
\cite[\S5.3]{OberdieckMC}. Extend it by fixing $1$ and the point
class. Acting factor by factor on Nakajima cycles defines
\begin{equation}\label{rdt:eq:Phi}
 \Phi_k\left(\prod_jq_{a_j}(\alpha_j)|0\rangle\right)
   =\prod_jq_{a_j}(\phi_k\alpha_j)|0\rangle.
\end{equation}
It is a graded isometric ring isomorphism and satisfies
$\Phi_k(\tau_i^{Y_d}(\alpha))=\tau_i^{S_k^{[d]}}(\phi_k\alpha)$.
The paper justifies the extension from parallel transport to real
isometries by Zariski density. This is a real cohomology comparison;
it is not automatically rational, integral, Hodge, or induced by a
functor on sheaves. In particular $D_k=\phi_kD$ can be nonalgebraic,
so its descendant is interpreted over $\R$ as above.

\begin{theorem}[Operator multiple-cover identity]
\label{rdt:thm:MC}
Over $\R$, one has the coefficientwise Laurent identity
\begin{equation}\label{rdt:eq:MC}
 \boxed{\displaystyle
 \widehat{\mathsf B}^D_{\beta,d}(r)=\sum_{k\mid\beta}
     \Phi_k^{-1}\widehat{\mathsf B}^{D_k}_{\beta_k,d}(r^k)\Phi_k.}
\end{equation}
There is no divisor weight multiplying the summand. In coefficients,
\begin{equation}\label{rdt:eq:MCcoeff}
 \widehat{\mathsf B}^D_{\beta,d,m}=\sum_{k\mid(\beta,m)}
     (-1)^{m-m/k}\Phi_k^{-1}
       \widehat{\mathsf B}^{D_k}_{\beta_k,d,m/k}\Phi_k.
\end{equation}
For the literal Euler-indexed tube,
\begin{equation}\label{rdt:eq:MCEu}
 \mathsf B^{D,\mathrm{Eu}}_{\beta,d}(r)
  =\sum_{k\mid\beta}r^{-d(k-1)}
       \Phi_k^{-1}\mathsf B^{D_k,\mathrm{Eu}}_{\beta_k,d}(r^k)\Phi_k.
\end{equation}
\end{theorem}
\begin{proof}
Apply the full relative-insertion identity
\cite[Theorem 5.1]{OberdieckMC} to the tube with the descendant
$\tau_0(\omega D)$ and arbitrary relative classes $a,b$. Its right
side uses $\phi_kD,\Phi_ka,\Phi_kb$, and substitutes $r^k$ for $r$.
By the isometry property, this says
\[
 \eta_d(\widehat{\mathsf B}^D_{\beta,d}a,b)
 =\sum_{k\mid\beta}
   \eta_d\bigl(\Phi_k^{-1}\widehat{\mathsf B}^{D_k}_{\beta_k,d}(r^k)\Phi_ka,b\bigr).
\]
Nondegeneracy gives \eqref{rdt:eq:MC}. Comparing $(-r)^m$ with
$(-r^k)^{m/k}$ gives the sign in \eqref{rdt:eq:MCcoeff}.
Substituting \eqref{rdt:eq:Eu} gives \eqref{rdt:eq:MCEu}.

The coefficient one here is a feature of the all-insertion theorem,
not an inference from a closed count. Its proof first reduces to caps
by the finite linear induction used in
Proposition~\ref{rdt:prop:Laurent}. For a cap, mixed fixed loci have
two cosections and hence vanish after only one reduction. Pure rubber
terms vanish in the required top insertion degree. A flag with $s$
nontrivial steps has a surjection of its obstruction sheaf to
$\mathcal O^s$; the flag differential maps to zero under the trace
maps. Removing one summand leaves a trivial obstruction for $s\ge2$,
so only a one-step locus remains. In rank one this is the explicit
argument of \cite[Proposition 4.8]{OberdieckMC}.

The one-step locus is a relative Quot scheme, equivalently the nested
Hilbert scheme of lengths $d_1,d$. Its localization factor contains
$(-1)^{(k-1)m/k}$ and powers of $kt$. Lemma 5.3 expresses every
stationary descendant as a universal polynomial in $kt$ and surface
intersection data. The dimension condition cancels the powers of
$kt$. The nested-Hilbert universality theorem
\cite[Theorem 3.2]{OberdieckMC} permits replacement of the data by
$\phi_k(\beta/k),\phi_k\alpha_i,\Phi_k\lambda_j$. Reversing the
calculation gives the displayed series substitution. Thus no factor
$1/k$ survives this top-degree calculation. The universality theorem
uses the nested virtual-cycle formula and Hilbert-scheme universality;
these are the geometric inputs of that argument.

In the printed equation (29), the substitution variable must be the
divisor of $\beta$: the proof ends with $p^r$ when its divisor is
called $r$. We have called that divisor $k$ to keep the series
variable $r$ distinct.
\end{proof}

This proves equality of cohomological endomorphisms, equivalently of
their transported Kunneth cohomology classes. It proves neither
equality of algebraic Chow correspondences under an unspecified
cycle map nor an equivalence of relative moduli or sheaf categories.

For a primitive effective charge in the marked lattice
\(N=U\oplus E_8(-1)\oplus E_7(-1)\), an isogeny supplies a
particular rational Hodge isometry.
Theorem~\ref{ist:thm:transport} extends its inverse normalized
pullback to the Hilbert states and proves the full insertion comparison.
It fixes the section and fibre and preserves the ordinary unit and pairing.
The signed twisted Mukai map of
Section~\ref{ist:sec:twist} has a different action:
it changes rank and does not preserve that ordinary pairing.
Its endpoint derived equivalence leaves the relative-operation map
of Problem~\ref{ist:prob:relative} to be constructed.

The same all-insertion theorem transports every reduced tensor used in
\eqref{rdt:eq:virtualglue}. The ordinary cup product, descendants and
diagonal are preserved by $\Phi_k$. Since each gluing has exactly one
reduced factor, its linear multiple-cover sum commutes with gluing.
This is a comparison of the first-order sewing systems. It is not a
claim that a sum of such transports is a multiplicative geometric
functor beyond this first-order recipient.

\begin{corollary}[The numerical divisor weight]\label{rdt:cor:closedMC}
Writing $Z_{\beta,d}(r)$ for the closed quotient series,
\begin{equation}\label{rdt:eq:closedMC}
 Z_{\beta,d}(r)=\sum_{k\mid\beta}\frac1k Z_{\beta_k,d}(r^k).
\end{equation}
In particular the imprimitive numerical multiple-cover formula is a
theorem, not an additional conjecture of this construction.
\end{corollary}
\begin{proof}
Take the trace of \eqref{rdt:eq:MC}. Theorem~\ref{rdt:thm:trace}
on the target K3 gives
\[
 \frac{\operatorname{Tr}\widehat{\mathsf B}^{D_k}_{\beta_k,d}}{\beta\cdot D}
 =\frac{\beta_k\cdot D_k}{\beta\cdot D}Z_{\beta_k,d}
 =\frac1k Z_{\beta_k,d}.
\]
The last equality is the surface isometry applied to
$(\beta/k,D)$. Thus the divisor weight is produced by the marked
quotient normalization upon closure.
\end{proof}

\subsection{Dependence on the marked class}

The correspondence is linear in $D$. Its decomposition
\eqref{rdt:eq:markeddecomposition} exhibits both its three-boundary
and cap dependence. Trace independence after dividing by $\beta D$
alone would not imply operator independence.

A stronger conclusion is available using an additional comparison.
Adopt the primary relative/Gromov--Witten identification quoted in
\cite[\S6.3]{OberdieckMC} from \cite[Corollary 4.1]{NesterovK3},
together with the Gromov--Witten divisor equation. The tautological
surface divisor $D^{[d]}$, for $d\ge2$, pairs with $\beta+mA$ as $\beta D$,
because it pairs trivially with the Hilbert--Chow extremal class $A$.
It follows that
\[
 \mu_{\beta,d,m}(D^{[d]},a)=(\beta D)\,A_{\beta,d,m}(a)
\]
for an operator $A$ independent of $D$.
For $d=1$ there is no extremal class; the three-relative contribution
has $m=0$, and the same surface divisor equation applies. For $d=0$,
the operation $\mu$ is zero by degree.

The cap term can be divided geometrically as well. With tangent
weight $-t$ at $0$, the one-step descendant is
$D^{[d]}-t(\beta D)$. After insertion of a point at the relative
evaluation, its first term vanishes. The localization calculation in
\cite[Theorem 4.2]{OberdieckMC} then yields
\begin{equation}\label{rdt:eq:capQ}
 \frac{c^D_{\beta,d,m}}{\beta D}
 =-\sum_{k\mid(\beta,m)}\frac{(-1)^m}{k}
    e^{\mathrm{vir}}\operatorname{Quot}\bigl(I_\eta,(0,\beta/k,m/k)\bigr).
\end{equation}
This is independent of $D$. Thus, with these stated extra inputs,
$\widehat{\mathsf B}^D/(\beta D)$ is independent of $D$ as an operator. This
conclusion uses the cap calculation and the primary divisor comparison;
it does not follow merely from the trace theorem.

\section{A section of an elliptic K3}
\label{rdt:sec:example}

Let $S$ have an elliptic fibration with section $C$ and fibre $F$.
Take $d=1$, $D=F$, and first $\beta=C$. Write
$p=[\mathrm{pt}]\in H^4(S)$, with $\int_Sp=1$. Define the
degree-zero operator
\begin{equation}\label{rdt:eq:T}
 T_C(a)=\left(\int_SaC\right)C.
\end{equation}
It vanishes on $H^0$ and $H^4$, acts nontrivially on $H^2$, and
satisfies $T_C^2=-2T_C$ and $\operatorname{Tr} T_C=-2$.

\subsection{A cap formula for every matrix coefficient}

For this subsection abbreviate the cap invariant at fixed $m$ by
$\operatorname{Cap}_m(I\mid a)$ and put
$\tau_b=\tau_0(\omega b)$. On the ordinary length-one cap,
$\tau_b$ gives the boundary class $b$. Degenerating a cap with two
interior insertions, putting $\tau_b$ on the new cap, gives
\begin{equation}\label{rdt:eq:capdifference}
 \eta_1(\widehat{\mathsf B}^D_{C,1,m}a,b)
 =\operatorname{Cap}_m(\tau_D\tau_b\mid a)
   -\operatorname{Cap}_m(\tau_b\mid Da).
\end{equation}
The second term is the ordinary-tube/reduced-cap summand. Omitting it
changes the operator.

Fix $t$ by $\mathrm{wt}(T_0\mathbb P^1)=-t$. Then
$[0]|_0=-t$ and the ideal $\mathcal O(-k[0])$ has weight $+kt$
at $0$. Its Chern character there is $e^{kt}$. In this convention,
directly from the fundamental horizontal point and vertical curve in
$\ch_2(\mathcal O_{\mathcal Z})$, one obtains on a one-step locus
\begin{equation}\label{rdt:eq:descweights}
 \tau_D=D_2-t(CD),\qquad
 \tau_b=b_2-t(Cb)\quad(b\in H^2(S)),\qquad
 \tau_1=1,\quad\tau_p=p_2.
\end{equation}
Here the subscript $2$ denotes the boundary length-one subscheme.
The exact sequence of $k[0]$ gives $1-e^{kt}$ for its equivariant
Chern character. An exponential $e^{-kt}$ instead uses the opposite
equivariant parameter. Keeping this convention fixed is necessary in
the local descendant calculation; polynomial universality and the
multiple-cover homogeneity argument are unchanged.

\begin{proposition}[The first two primitive kernels]
\label{rdt:prop:firstkernels}
For $\beta=C$, $d=1$, $D=F$,
\begin{align}
 K^F_{C,1,0}&=C\boxtimes C,&\widehat{\mathsf B}^F_{C,1,0}&=T_C,
 \label{rdt:eq:K0}\\
 K^F_{C,1,1}&=[\Delta_S],&\widehat{\mathsf B}^F_{C,1,1}&=\id_{H^*(S)}.
 \label{rdt:eq:K1}
\end{align}
These two cases correspond to tube Euler characteristics $N=1,2$,
respectively.
The equalities of kernels here are in cohomology.
\end{proposition}
\begin{proof}
The cap reduction and one-step localization are those of
\cite[\S\S4.4--5.3]{OberdieckMC}. For primitive $C$ the one-step
length is $d_1=m$.

For $m=0$ the nested Hilbert scheme is $S_C^{[0,1]}=C$: the
inclusion $\mathcal O(-C)\subset I_x$ holds precisely when $x\in C$.
Its reduced class is $[C]$. The virtual normal Euler factor begins
with $t^{-2}$. For $a,b\in H^2(S)$ put $A=aC$, $B=bC$, and write
$h\in H^2(C)$ with $\int_Ch=1$. The relevant restrictions are
$a_2=Ah$, $b_2=Bh$, $D_2=h$. The first term of
\eqref{rdt:eq:capdifference} is
\[
 \int_C t^{-2}(h-t)(Bh-tB)Ah=AB.
\]
Only the $t^2h$ term survives. Higher terms of the normal Euler
factor cannot contribute in this dimension. The second cap term
vanishes because $Da$ restricts to zero on $C$. The pairings
$(1,p)$ and $(p,1)$ vanish because $p|_C=0$; all other degree
combinations vanish by virtual dimension. Hence the full bilinear
form is $(aC)(bC)$, proving \eqref{rdt:eq:K0}. Geometrically, its
moving degree-one graph maps have both boundary points on $C$.

For $m=1$, consider $S_C^{[1,1]}$ in $S_1\times S_2$. Put
$L_i=\operatorname{pr}_i^*\mathcal O(C)$ and $C_i=c_1(L_i)$.
The derived surface Hom calculation gives
\[
 R\operatorname{Hom}(I_1,I_2(C))
       =\mathcal O-L_1-L_2+\mathcal O_\Delta L_1.
\]
For example this follows by expanding $I_i=\mathcal O-\mathcal O_{x_i}$;
the derived self-duality of a point on the K3 has even shift two,
and the coincident-point term is $\mathcal O_\Delta L_1$.
Since $R\Gamma\mathcal O(C)=\mathcal O$ and $|C|$ is a point,
the nested virtual-cycle formula
\cite[Theorem 3.1]{OberdieckMC} gives
\begin{equation}\label{rdt:eq:W}
 W_2:=j_*[S_C^{[1,1]}]^\mathrm{red}
     =c_2(L_1+L_2-\mathcal O_\Delta L_1)
     =C_1C_2+[\Delta_S].
\end{equation}
Indeed $\ch_2(\mathcal O_\Delta L_1)=[\Delta_S]$, which verifies
the last Chern-class identity directly.

Let $K=I_1(-C)$ and $Q=I_2-K$ in the universal quotient sequence.
The surface Hom class controlling the moving Euler factor is
\[
 V=R\operatorname{Hom}(K,Q)
     =T_{S_1}-\mathcal O-L_1-L_2+\mathcal O_\Delta L_1.
\]
It has rank $-1$ and $c_1(V^\vee)=C_1+C_2$. Thus
\begin{equation}\label{rdt:eq:eulerfactor}
 e(V^\vee\otimes\mathbf t^{-1})
   =-t^{-1}+(C_1+C_2)t^{-2}+O(t^{-3}).
\end{equation}
For $a,b\in H^2(S)$, the constant equivariant coefficient of the
first cap integral is
\[
 \int_{S\times S}W_2\left\{
  a_2\bigl((CD)b_2+(Cb)D_2\bigr)
    +(CD)(Cb)a_2(C_1+C_2)\right\}.
\]
The last summand integrates to zero, because
\[
 \int W_2a_2C_1=(C^2+1)(aC)=-(aC),\qquad
 \int W_2a_2C_2=aC.
\]
The first cap therefore equals $(CD)(ab)+(Cb)(aD)$, where
juxtaposed degree-two classes in parentheses denote their surface
intersection. The second cap in \eqref{rdt:eq:capdifference}
equals $(Cb)(aD)$. Subtraction gives $(CD)(ab)=\eta_1(a,b)$.

For $(a,b)=(1,p)$, the sole contribution is
$(-tCD)p_2(-t^{-1})$ integrated against $W_2$, giving $CD=1$;
the second cap is zero. The same holds for $(p,1)$. Every other
graded pairing vanishes by dimension. This is the full Poincare
bilinear form on $H^0\oplus H^2\oplus H^4$, proving
\eqref{rdt:eq:K1} by Kunneth and Poincare nondegeneracy.
\end{proof}

The $+1$ on the unit and top class in \eqref{rdt:eq:K1} is also
visible on the cap over a generic boundary point away from $C$.
There $\operatorname{Quot}(I_x,(0,C,1))$ is a reduced virtual point,
$R\operatorname{Hom}(K,Q)=-\mathcal O$, and the two equivariant
factors are $-1/t$ and $-t(CD)$. Their product is $+1$.

\subsection{The full pairing and the imprimitive class}

Use a basis
\[
 1,C,F,u_2,v_2,u_3,v_3,e_1,\ldots,e_8,f_1,\ldots,f_8,p
\]
of $H^*(S,\Q)$. The $H^2$ Gram matrix is
\begin{equation}\label{rdt:eq:K3pair}
 \begin{pmatrix}-2&1\\1&0\end{pmatrix}
 \oplus U\oplus U\oplus E_8(-1)\oplus E_8(-1),
 \qquad U=\begin{pmatrix}0&1\\1&0\end{pmatrix}.
\end{equation}
The $H^0,H^4$ pairing is $(1,p)=(p,1)=1$, with all other
cross-degree blocks zero. This is the full K3 Poincare pairing.
The only nonzero row of $T_C$ is its $C$ row, with entries $-2,1$
in columns $C,F$. Its nonzero bilinear block is
\[
 \bigl(\eta_1(T_Ca,b)\bigr)_{a,b=C,F}
  =\begin{pmatrix}4&-2\\-2&1\end{pmatrix}.
\]
Contraction with the full inverse of \eqref{rdt:eq:K3pair}, including
the $1,p$ block, gives $\operatorname{Tr} T_C=-2$ and $\operatorname{Tr}\id=24$.

\begin{proposition}[The doubled section]\label{rdt:prop:double}
For $\beta=2C$, $d=1$ and $D=F$,
\begin{equation}\label{rdt:eq:double}
 \widehat{\mathsf B}^F_{2C,1}(r)=\widehat{\mathsf B}^F_{C,1}(r^2),
\end{equation}
and, in particular,
\[
 \widehat{\mathsf B}^F_{2C,1,0}=T_C,\qquad
 \widehat{\mathsf B}^F_{2C,1,1}=0,\qquad
 \widehat{\mathsf B}^F_{2C,1,2}=-\id.
\]
\end{proposition}
\begin{proof}
The negative intersection with $C$ implies inductively that
$|jC|=\{jC\}$. Riemann--Roch gives $h^1(\mathcal O(2C))=3$:
$h^0=1$, $h^2=0$ and $\chi=-2$. In the one-step contribution
with divisor $k=1$, the ambient pushforward of the nested virtual
class contains the factor
$c_1(\mathcal O_{|2C|}(1))^3=0$ by
\cite[Theorem 3.1]{OberdieckMC}. All the cap integrands are pullbacks
of universal classes on that ambient product, so these contributions
vanish. This conclusion concerns the required virtual integrals; no
intrinsic-cycle vanishing stronger than the pushforward formula is
needed.

For $k=2$, use $S_2=S$ and $\phi_2=\id$, since $(2C)/2=C$ is
primitive. The one-step calculation, followed by the linear cap
induction, is exactly the $C$ calculation with $r$ replaced by $r^2$.
This proves \eqref{rdt:eq:double}. Extracting coefficients with the
sign $(-r)^m$ gives a factor $(-1)^{m-m/2}$ when $m$ is even;
there is no such contribution when $m$ is odd. Proposition
\ref{rdt:prop:firstkernels} now gives the three displayed coefficients.
\end{proof}

The vanishing at $m=1$, which is relative Euler
characteristic $N=2$, has a geometric source. The $k=1$ contribution
has vanishing ambient reduced virtual class, and the $k=2$ thickening
cannot have odd third Chern character. The next coefficient is
$N=3$, with the nonzero operator $-\id$ in the same family.

\subsection{One further primitive coefficient}

With the primary/Gromov--Witten comparison stated at the end of
Section~\ref{rdt:sec:covers}, the length-one three-relative operation
is concentrated in $m=0$. An honest graph of a genus-zero tree of
maps to $S$ has Euler characteristic one and degree one over the
parametrized base, hence $m=0$. There is no Hilbert--Chow curve in
$S^{[1]}=S$ giving another index.

For the cap, put
$Q_{1,C,m}=e^{\mathrm{vir}}\operatorname{Quot}(I_x,(0,C,m))$. The leading
$q^{-1}$ coefficient of \cite[Theorem 1.2]{OberdieckMC} gives
\[
 \sum_m Q_{1,C,m}r^m=\frac{r}{(1-r)^2}.
\]
Indeed the leading theta square is $(r-1)^2/r$ and the leading
coefficient of the function $G$ in that theorem is one. Formula
\eqref{rdt:eq:capQ} thus gives $c^F_{C,1,m}=(-1)^{m-1}m$ for
$m\ge1$. Consequently, with this additional primary comparison,
\begin{equation}\label{rdt:eq:fullC}
 \widehat{\mathsf B}^F_{C,1}(r)=T_C-\frac{r}{(1-r)^2}\id,
 \qquad \widehat{\mathsf B}^F_{C,1,2}=-2\id.
\end{equation}
The direct computations of $m=0,1$ in Proposition
\ref{rdt:prop:firstkernels}, and the doubled coefficient at $m=2$,
do not require this extra identification.

Both index conventions and all relevant contractions are displayed
together below. The final column is the coefficient of $(-r)^m$
in the closed series.
\begin{center}
\begin{tabular}{ccccrr}
\toprule
Class & tube $N$ & closed $m$ & operator & trace & closed invariant\\
\midrule
$C$  & 1 & 0 & $T_C$ & $-2$ & $-2$\\
$C$  & 2 & 1 & $\id$ & $24$ & $24$\\
$C$  & 3 & 2 & $-2\id$ & $-48$ & $-48$\\
$2C$ & 1 & 0 & $T_C$ & $-2$ & $-1$\\
$2C$ & 2 & 1 & $0$ & $0$ & $0$\\
$2C$ & 3 & 2 & $-\id$ & $-24$ & $-12$\\
\bottomrule
\end{tabular}
\end{center}

\subsection{Two incompatible sewing rules}

Suppose one were to put the closed numerical divisor weight into the
marked tube, replacing the $k=2$ term in
\eqref{rdt:eq:double} by $\tfrac12\widehat{\mathsf B}^F_{C,1}(r^2)$.
At the first nonzero term, $m=0$, closure would give
\[
 \frac{\operatorname{Tr}(\tfrac12T_C)}{2C\cdot F}=-\frac12,
 \qquad\text{whereas}\qquad
 \frac{\operatorname{Tr} T_C}{2C\cdot F}=-1.
\]
At the first positive $k=2$ term, $m=2$, it would give
\begin{equation}\label{rdt:eq:badweight}
 \frac{\operatorname{Tr}(-\tfrac12\id_{24})}{2}=-6,
 \qquad\text{whereas}\qquad
 \frac{\operatorname{Tr}(-\id_{24})}{2}=-12.
\end{equation}
This difference is an exact obstruction to that rule. It already
uses the full state space, rather than only its middle cohomology.

There is a separate obstruction to deleting the ordinary factor from
reduced degeneration. An unmarked collar on the marked tube satisfies
\begin{equation}\label{rdt:eq:collar}
 \widehat{\mathsf B}^D=\widehat{\mathsf B}^D\id+M_{D^{[d]}}R_{\beta,d}
         =\widehat{\mathsf B}^D\id,
 \qquad R_{\beta,d}=0.
\end{equation}
Replacing the ordinary collar by the unmarked reduced collar gives
$\widehat{\mathsf B}^D R_{\beta,d}=0$. At $(C,d,m)=(C,1,0)$ this would assert
$T_C=0$, although $T_C(F)=C$ and $T_C(C)=-2C$.
Equivalently, retaining only reduced factors would put the product
in $\epsilon^2=0$. If one instead ignores the first-order grading
and multiplies two marked matrices, $T_C^2=-2T_C\ne T_C$ still fails
the collar equation. None of these replacements respects the unit.


\section{An intrinsic first-order class}
\label{rdt:sec:Hochschild}

The section kernel detects more than a failure of one formula for
sewing. It detects a deformation of the cup algebra that cannot be
removed by changing its first-order splitting. Retain the section
$C$, fibre $F$ and point class $p$ of the preceding example, so
\[
 C^2=-2p,\qquad F^2=0,\qquad CF=p,\qquad \int_Sp=1.
\]
Put $A=H^*(S,\Q)$ and $B=\Q[\epsilon]/(\epsilon^2)$, with
$|\epsilon|=2$. At the single coefficient $(\beta,d,m)=(C,1,0)$,
write $\mu=\mu_{C,1,0}$. Proposition~\ref{rdt:prop:firstorder}
makes
\[
 a\star b=ab+\epsilon\mu(a,b)
\]
a unital graded associative product on $A\otimes_\Q B$.
Here $ab$ is cup product and $\mu$ has internal degree $-2$.

For a linear one-cochain $f:A\to A$, the Hochschild coboundary is
\[
 (\delta f)(a,b)=af(b)-f(ab)+f(a)b.
\]
The cocycle equation for a two-cochain is
\[
 a\mu(b,c)-\mu(ab,c)+\mu(a,bc)-\mu(a,b)c=0.
\]
It is precisely \eqref{rdt:eq:assoc}. All classes of $A$ are even,
so these formulas have no additional parity signs. Hochschild
cohomology $HH^2(A,A)$ is the space of two-cocycles modulo these
coboundaries, with its additional internal grading.

\begin{proposition}[The section deformation]
\label{rdt:prop:HHnonzero}
The class of $\mu$ is nonzero in the internal degree $-2$ part of
$HH^2(A,A)$. The algebra $(A\otimes B,\star)$ is not isomorphic
as a $B$-algebra to the undeformed cup algebra $A\otimes B$,
even if the reduction of the isomorphism modulo $\epsilon$ is
allowed to be an arbitrary algebra automorphism of $A$.
\end{proposition}
\begin{proof}
At the chosen coefficient, Proposition~\ref{rdt:prop:firstkernels}
gives $\widehat{\mathsf B}^{F}_{C,1,0}=T_C$. The marked-state
identity \eqref{rdt:eq:markeddecomposition} says
\[
 T_C(a)=\mu(F,a)+c^F_{C,1,0}a.
\]
Since $T_C(1)=0$ and $\mu(F,1)=0$, the scalar $c^F_{C,1,0}$ is
zero. In particular,
\begin{equation}\label{rdt:eq:HHcoefficient}
 \mu(F,F)=T_C(F)=C.
\end{equation}

Define a functional on two-cochains by
\[
 \mathcal L(\nu)=\int_S F\nu(F,F).
\]
For every linear $f$, not necessarily homogeneous, $F^2=0$ gives
\[
 (\delta f)(F,F)=2Ff(F),\qquad
 \mathcal L(\delta f)=2\int_SF^2f(F)=0.
\]
On the other hand, \eqref{rdt:eq:HHcoefficient} gives
$\mathcal L(\mu)=\int_SFC=1$. Thus $\mu$ is not a coboundary.

An isomorphism to the undeformed algebra that reduces to the
identity has the form $T(a)=a+\epsilon f(a)$. Comparing
$T(a\star b)$ with $T(a)T(b)$ gives $\mu=\delta f$, a
contradiction. If the reduction is an automorphism $g$ instead,
compose $T$ with the constant extension of $g^{-1}$ in the
undeformed target. The resulting isomorphism reduces to the
identity, so the same contradiction applies.
\end{proof}

This obstruction concerns the full ordinary cohomology algebra.
It is already detected by the square-zero fibre class and the
rank-one correspondence $C\boxtimes C$. No scalar trace quotient
or choice of a complement to $\langle C,F\rangle$ is involved.

The consequence for a comparison of operations is explicit.
If two first-order algebras have products
$m_A+\epsilon\nu_A$ and $m_H+\epsilon\nu_H$, and a $B$-linear
map is written $J=J_0+\epsilon J_1$, its first-order
multiplicativity equation is
\begin{equation}\label{rdt:eq:deformedcomparison}
\begin{aligned}
 J_1(m_A(a,b))+J_0(\nu_A(a,b))
  ={}&m_H(J_1(a),J_0(b))+m_H(J_0(a),J_1(b))\\
    &+\nu_H(J_0(a),J_0(b)).
\end{aligned}
\end{equation}
Indeed, these are the two coefficients of
$J(a\star_A b)=J(a)\star_H J(b)$ after the ordinary
multiplicativity of $J_0$ is imposed. In the relative sewing
algebra, multiplication of general inputs is
\begin{equation}\label{rdt:eq:allfirstorderterms}
 (a_0+\epsilon a_1)\star(b_0+\epsilon b_1)
 =a_0b_0+\epsilon\bigl(a_1b_0+a_0b_1+\mu(a_0,b_0)\bigr).
\end{equation}
The last term is an operation on two ordinary inputs. At $F,F$
it is $C$, and Proposition~\ref{rdt:prop:HHnonzero} shows that a
change of splitting cannot remove it. The fixed-multiplication
lemma~\ref{ff:lem:dual-numbers} remains valid: its source and
target multiplications do not vary with $\epsilon$. Applying that
lemma to a geometrically deformed recipient requires the extra
terms in \eqref{rdt:eq:deformedcomparison}.

\subsection{The complete section coefficient}

The primary three-boundary invariant determines the entire product
$\mu_{C,1,0}$, including its values orthogonal to the section and fibre.
Write
\[
 A=\Q\,1\oplus V\oplus\Q\,p,\qquad V=H^2(S,\Q),
 \qquad q(v,w)=\int_Svw.
\]
Thus $vw=q(v,w)p$ for $v,w\in V$, while $pV=p^2=0$.
Define the linear functional
\[
 s_C(a)=\int_SaC,\qquad a\in A.
\]
It vanishes on $H^0(S,\Q)$ and $H^4(S,\Q)$.

\begin{proposition}[The section cubic]
\label{rdt:prop:sectioncubic}
For every $a,b\in H^*(S,\Q)$,
\[
 \mu_{C,1,0}(a,b)=s_C(a)s_C(b)C.
\]
Equivalently, its paired trilinear form is
\begin{equation}\label{rdt:eq:sectioncubic}
 \eta_1\bigl(\mu_{C,1,0}(a,b),c\bigr)
       =s_C(a)s_C(b)s_C(c).
\end{equation}
In particular,
\[
 \mu(C,C)=4C,\qquad \mu(C,F)=-2C,\qquad \mu(F,F)=C.
\]
Every pair involving a class orthogonal to $C$ has zero product.
\end{proposition}

\begin{proof}
The primary relative Donaldson--Thomas/Gromov--Witten comparison
\cite[Definition 3.3 and Corollary 4.1]{NesterovK3}
identifies the paired coefficient with the primary reduced
genus-zero invariant of $S^{[1]}=S$ in class $C$.
Its insertions range over the full rational cohomology.

For $v\in H^2(S,\Q)$, the primary divisor equation contributes
$\int_Cv$. Indeed, the evaluation pullback integrates to this number
on the fibre of the forgetful universal curve.
There is no descendant correction for primary insertions.
Applying this identity to three degree-two classes gives
\[
 \eta_1\bigl(\mu(a,b),c\bigr)
      =N_C\,s_C(a)s_C(b)s_C(c),\qquad a,b,c\in V,
\]
where $N_C$ is the zero-point reduced invariant.
The geometric section kernel fixes its normalization.
Equation~\eqref{rdt:eq:HHcoefficient} and $s_C(F)=1$ give
\[
 N_C=\eta_1\bigl(\mu(F,F),F\bigr)=\int_SCF=1.
\]

The trilinear coefficient has total real cohomological degree six.
The possible even-degree triples are $(2,2,2)$ and permutations
of $(0,2,4)$. The latter vanish by the unit identity and cyclicity.
This proves \eqref{rdt:eq:sectioncubic} on all of $A$.
Nondegeneracy of the full Poincare pairing, including its $1,p$
block, gives the asserted product.
\end{proof}

\subsection{A cubic undetected by the marked operator and traces}

The fibre-marked operator and the multiplication traces do not
determine an arbitrary cyclic deformation of this cohomology algebra.
For $w\in V$, put
\[
 \ell_w(a)=\int_Saw,\qquad
 \nu_w(a,b)=\ell_w(a)\ell_w(b)w.
\]
The map $\nu_w$ vanishes when either input has no degree-two component.

\begin{lemma}\label{rdt:lem:rankonecubic}
The cochain $\nu_w$ is a normalized cyclic Hochschild two-cocycle
of internal degree $-2$.
For a parameter $t$ of degree two, the product
\[
 a*_w b=ab+t\nu_w(a,b)
\]
is associative and unital over $\Q[t]$.
Its invariant pairing is the original Poincare pairing.
\end{lemma}

\begin{proof}
Let $a_0$ denote the coefficient of $1$ in $a$.
The cup algebra gives
\[
 \ell_w(ab)=a_0\ell_w(b)+b_0\ell_w(a),
 \qquad
 aw=a_0w+\ell_w(a)p.
\]
Substitution into the Hochschild differential yields
\[
 a\nu_w(b,c)-\nu_w(ab,c)+\nu_w(a,bc)-\nu_w(a,b)c=0.
\]
The terms proportional to $w$ cancel in pairs.
The two terms proportional to $p$ are
$\ell_w(a)\ell_w(b)\ell_w(c)p$ with opposite signs.
This verifies the identity for arbitrary states of $A$.

The coefficient of $t^2$ in associativity also vanishes, since
\[
 \nu_w(\nu_w(a,b),c)
 =q(w,w)\ell_w(a)\ell_w(b)\ell_w(c)w
 =\nu_w(a,\nu_w(b,c)).
\]
The unit is unchanged because $\ell_w(1)=0$.
Finally,
\[
 \eta_1(\nu_w(a,b),c)=\ell_w(a)\ell_w(b)\ell_w(c)
\]
is symmetric. Thus the original pairing is invariant.
\end{proof}

Choose $u=u_2$ and $v=v_2$ in the second hyperbolic plane of
\eqref{rdt:eq:K3pair}. These classes satisfy
\[
 q(u,u)=q(v,v)=0,\qquad q(u,v)=1,\qquad
 u,v\perp C,F.
\]
Set
\[
 \widetilde\mu=\mu_{C,1,0}+\nu_u.
\]

\begin{proposition}\label{rdt:prop:traceinvisiblecubic}
The products $ab+\epsilon\mu(a,b)$ and
$ab+\epsilon\widetilde\mu(a,b)$ are unital cyclic deformations of
the full cohomology algebra over $\Q[\epsilon]/(\epsilon^2)$.
They have the same fibre-marked operator.
For every $D\in A$, their multiplication-correction operators
have the same trace:
\[
 \operatorname{Tr}\bigl(\widetilde\mu(D,-)\bigr)
   =\operatorname{Tr}\bigl(\mu(D,-)\bigr).
\]
Their Hochschild classes are distinct when the ordinary
identification of $A$ is fixed.
\end{proposition}

\begin{proof}
Proposition~\ref{rdt:prop:sectioncubic} identifies $\mu$ with $\nu_C$.
The mixed quadratic compositions of $\nu_C$ and $\nu_u$ vanish
because $q(C,u)=0$.
Lemma~\ref{rdt:lem:rankonecubic} therefore also gives an associative
polynomial product $ab+t\widetilde\mu(a,b)$.
Reducing modulo $t^2$ gives the stated first-order deformation.

Since $\ell_u(F)=0$,
\[
 \widetilde\mu(F,a)=\mu(F,a)\qquad(a\in A).
\]
The cap scalar in the fibre-marked operator remains zero.
For any $D$, the difference of the multiplication operators is
the rank-one endomorphism
\[
 \nu_u(D,-)=\ell_u(D)\,u\otimes\ell_u.
\]
Its trace on all twenty-four states is
$\ell_u(D)\ell_u(u)=0$.
Thus the complete fibre-marked operator and all the stated traces
agree. Holding the cap scalars fixed also preserves the
marked-tube traces.

The untraced product distinguishes the two deformations:
\[
 \mu(v,v)=0,\qquad
 \widetilde\mu(v,v)=u,\qquad
 \int_Sv\widetilde\mu(v,v)=1.
\]
Define $\mathcal L_v(\theta)=\int_Sv\theta(v,v)$.
For every linear one-cochain $f$, the equality $v^2=0$ gives
\[
 \mathcal L_v(\delta f)
       =2\int_Sv^2f(v)=0.
\]
But $\mathcal L_v(\nu_u)=1$.
Hence $\nu_u$ is not a Hochschild coboundary, and the two
first-order classes cannot agree through a change of splitting
that fixes the ordinary identification.
\end{proof}

The polynomial deformation in this construction is algebraic.
It is not identified with higher-order geometric reduced
Donaldson--Thomas coefficients.
An operation-preserving geometric comparison must recover
\eqref{rdt:eq:sectioncubic}, including its zero value at $(v,v,v)$.
Agreement of the marked operator and all marked traces leaves
this coefficient undetermined.

\section{The primitive inverse Igusa series}
\label{rdt:sec:Igusa}

For a primitive effective class $\beta_h$ with square $2h-2$, choose
$D_h$ with $\beta_hD_h\ne0$. The geometric trace supplies the
closed DT expansion
\[
 Z_{h,d}^{0}(r)=\frac{1}{\beta_hD_h}\operatorname{Tr}\widehat{\mathsf B}^{D_h}_{\beta_h,d}(r)
\]
at $r=0$. The primitive evaluation quoted in
\cite[Remark 1.2, equation (3)]{OberdieckMC} is the inverse Igusa
theorem of \cite{OB}. The comparison of reduced DT and
PT series and rationality is the theorem of
\cite{OberdieckReducedPT}.

To distinguish expansions, let $R_{h,d}(v)$ be the rational function
whose Laurent expansion at $v=0$ is the reduced DT/PT series.
The variables are
\begin{equation}\label{rdt:eq:OPvars}
 (T_{\mathrm{OP}},Q_{\mathrm{OP}},P_{\mathrm{OP}})=(s,q,r),
 \qquad v=-r.
\end{equation}
Thus $Z_{h,d}^{0}(r)$ is the expansion of $R_{h,d}(-r)$ at zero.
The rational representative of the complete primitive trace is
\begin{equation}\label{rdt:eq:Igusa}
 \sum_{h,d\ge0}R_{h,d}(-r)q^{h-1}s^{d-1}
   =-\frac1{\chi_{10}^{\mathrm{OP}}(r,q,s)}.
\end{equation}
In particular the boundary-length exponent is $d-1$, not $d$.

In the theta convention
\[
 \Delta_5=64(qrs)^{1/2}
  \prod_{(h,\ell,d)>0}(1-q^hr^\ell s^d)^{f(hd,\ell)},
 \qquad D_5=\Delta_5/64,
\]
and the monic Igusa product has exponents $2f$, so
\begin{equation}\label{rdt:eq:Igusanorm}
 \chi_{10}^{\mathrm{OP}}=D_5^2=4096^{-1}\Delta_5^2,
 \qquad Z^X_{\mathrm{OP}}=-4096\Delta_5^{-2}.
\end{equation}
Thus the normalized geometric trace in \eqref{rdt:eq:trace}, summed
with $q^{h-1}s^{d-1}$, realizes $Z^X_{\mathrm{OP}}$.
Multiplication of that scalar functional by $-1/4096$ realizes
$\Delta_5^{-2}$. This rescaling is not a definition of a protected
vacuum functional.

For fixed $h,d$ the product coefficient is rational in $r$: apart
from the leading factor from $(h,d)=(0,0)$, only finitely many
factors can contribute at fixed nonnegative $q,s$ degree. The weak
Jacobi bound leaves finitely many $\ell$ at each such degree. The
leading inverse factor is $r/(1-r)^2$. The rational comparison in
\eqref{rdt:eq:Igusa} identifies this representative with the actual
DT expansion. Expansion at $r=\infty$ in the product chamber and
expansion at $r=0$ for DT are distinct operations on this rational
function. The substitution $r=e^{iu}$ for Gromov--Witten theory
is a third expansion. None of these operations replaces another
without the rational-representative theorem.

The term $d=0$ is indispensable. For the section class its primitive
coefficient is
\[
 [q^{-1}s^{-1}]Z^X_{\mathrm{OP}}
    =-\frac{r}{(1-r)^2}.
\]
It is furnished by the scalar endomorphism on $H_0=\Q$.
Deleting the empty relative boundary would delete this entire
$s^{-1}$ coefficient.

Finally, let
$c(a,b,j)=[q^as^br^j](\chi_{10}^{\mathrm{OP}})^{-1}$ in the
specified coefficient expansion. Corollary~\ref{rdt:cor:closedMC}
gives
\[
 \mathrm{DT}^{S\times E}_{n,(\beta,d)}
 =(-1)^{n+1}\sum_{k\mid(\beta,n)}\frac1k
        c\left(\frac{(\beta/k)^2}{2},d-1,n/k\right),
\]
the already established numerical theorem
\cite[Theorem 1.3]{OberdieckMC}. It is the scalar consequence of
the preceding all-boundary comparison.

\section{The length-two boundary state}
\label{rdt:sec:Hilb2}

The next boundary space has two Nakajima partitions. Keeping both
is necessary for its diagonal. The following description fixes all
normalizations independently of a DT kernel.

Let $B=\operatorname{Bl}_\Delta(S\times S)$, with blowdown $\pi$ and degree-two
quotient $f:B\to Y_2=S^{[2]}$. Let
$j:E_{\mathrm{exc}}\hookrightarrow Y_2$ and
$p_E:E_{\mathrm{exc}}=\mathbb P(TS)\to S$. Fix the positive
geometric creation convention
\[
 U(a,b)=f_*\pi^*(a\boxtimes b)=q_1(a)q_1(b)|0\rangle,
 \qquad V(a)=j_*p_E^*a=q_2(a)|0\rangle.
\]
There is no factorial in $U$, and $V(1)$ is the effective exceptional
divisor. Their degrees are $|U(a,b)|=|a|+|b|$ and
$|V(a)|=|a|+2$.

\begin{proposition}[The full length-two pairing]
\label{rdt:prop:Hilb2}
The state space is
\[
 H^*(Y_2,\Q)\cong\Sym^2H^*(S,\Q)\oplus H^*(S,\Q)[-2],
\]
with the degree convention just stated, and
\begin{align}
 (U(a,b),U(c,d))_2&=(a,c)_1(b,d)_1+(a,d)_1(b,c)_1,
 \label{rdt:eq:UU}\\
 (V(a),V(b))_2&=-2(a,b)_1,\label{rdt:eq:VV}\\
 (U(a,b),V(c))_2&=0.\label{rdt:eq:UV}
\end{align}
For $e_i$ a homogeneous basis and $e^i$ its Poincare dual, a basis is
$U(e_i,e_j)$ for $i\le j$ together with $V(e_i)$. Its dual basis is
\[
 \frac{U(e^i,e^j)}{1+\delta_{ij}},\qquad -\frac12V(e^i).
\]
The cup unit is $U(1,1)/2$ and the normalized top class is $U(p,p)$.
\end{proposition}
\begin{proof}
The blowup decomposition contributes the diagonal exceptional copy
of $H^*(S)$ shifted up by two real degrees. Taking invariants for
the transposition on the degree-two quotient gives the symmetric
square of the remaining $H^*(S)\otimes H^*(S)$. This proves the
description of the state and its two sectors.

The degree-two quotient gives
\[
 f^*U(c,d)=\pi^*(c\boxtimes d+d\boxtimes c).
\]
The projection formula now proves \eqref{rdt:eq:UU}. If $E_B$ is
the exceptional divisor of the blowup, then $f^*E_{\mathrm{exc}}=2E_B$
and $f$ restricts to an isomorphism on the exceptional divisor.
Since $N_{E_B/B}=\mathcal O(-1)$, this gives
$N_{E_{\mathrm{exc}}/Y_2}=\mathcal O(-2)$. The self-intersection
formula and integration over its projective-line fibres give
\eqref{rdt:eq:VV}. Restricting $U(a,b)$ to that divisor gives
$2p_E^*(ab)$; it contains no fibre hyperplane class, so its integral
against $p_E^*c$ is zero. This proves \eqref{rdt:eq:UV}.

The dual-basis formulas follow by substituting into these three
pairings. The factors two for repeated symmetric basis elements and
$-2$ for the exceptional sector are essential. Finally $f_*1=2$
and $\int_{Y_2}U(p,p)=1$, giving the asserted unit and top class.
\end{proof}

There are $24\cdot25/2=300$ symmetric states and $24$ exceptional
states, for a total of $324$. The Betti numbers in real degrees
$0,2,4,6,8$ are $(1,23,276,23,1)$. The complete coevaluation is
\[
 \sum_{i\le j}U(e_i,e_j)\otimes
          \frac{U(e^i,e^j)}{1+\delta_{ij}}
   -\frac12\sum_iV(e_i)\otimes V(e^i).
\]
Its contraction is $300+24=324$.

An insertion-sensitive length-two evaluation of the relative DT
kernel must compute its entries on both sectors and contract with
this diagonal. For example the multiple-cover theorem fixes in
advance the $k=2$ comparison, including the coefficient sign
$(-1)^{m-m/2}$ and the factor $1/2$ only after closure. The state
and pairing above do not themselves evaluate those DT entries.

\section{The critical coefficient comparison}
\label{rdt:sec:critical}

Relative DT sewing supplies a geometric cohomological recipient and
its closed numerical map. It is useful because the full boundary
amplitudes contain information that the inverse Igusa scalar does
not determine. The rank-one kernel $C\boxtimes C$, the diagonal
kernel, and the ordinary collar have different actions on the unit,
middle cohomology and top class, even before a trace is considered.
The construction requires neither a BKM recognition theorem nor a
choice of a prepared critical state.

Another geometric carrier, the Hilbert-scheme elliptic genus, retains
its own Fock state and yields an Igusa product with a boundary factor.
An equality of such indices cannot identify its boundary operations
with the relative ideal-sheaf correspondences. The latter keep the
K3 curve class, stationary descendant, reduced virtual theory and
multiple-cover comparison as specified data.

A compact critical Hall realization has a further first requirement.
For an extension correspondence
\[
 \mathcal Q_\alpha\times_{BE}\mathcal Q_\beta
   \xleftarrow{\ p\ }\mathcal E_{\alpha,\beta}
   \xrightarrow{\ q\ }\mathcal Q_{\alpha+\beta},
\]
one needs independently defined reduced oriented critical coefficient
complexes $\mathscr K_\gamma$ and a map
\begin{equation}\label{rdt:eq:criticalmap}
 p^*(\mathscr K_\alpha\boxtimes_{BE}\mathscr K_\beta)
       [\kappa_{\alpha,\beta}]
       \longrightarrow q^!\mathscr K_{\alpha+\beta}.
\end{equation}
The chosen support must make $q$ proper for the pushforward entering
the trace. The orientation square roots, their comparison maps,
the shifts and base-change maps must be compatible on every
extension flag, including nonsplit three-step flags and their
associativity homotopies. The common-$BE$ zero-object unit must
act as the identity under these same comparisons.
The common $E$-torsor in this fibre product is part of the extension
geometry. It is different from the quotient-stack normalization in
Theorem~\ref{rdt:thm:trace}.

\begin{question}\label{rdt:qn:realization}
Given the critical coefficient system and its flag-compatible maps,
construct a boundary realization to $H_d$ that intertwines its marked
state action with \eqref{rdt:eq:kernel}, its ordinary and reduced
sewings with \eqref{rdt:eq:virtualglue}, and its supported trace with
\eqref{rdt:eq:trace}. Compare the independent state parameter family,
line and section with the modular line and section, including their
characters and boundary normalizations.
\end{question}

Any comparison with this relative sewing system must also preserve
the intrinsic first-order class of
Proposition~\ref{rdt:prop:HHnonzero}, through its actual coefficient
map. In particular, a $B$-linear identification of the target with
an undeformed cup algebra cannot supply the comparison.

The equality of all ordinary cohomological matrix coefficients under
$\Phi_k$ does not construct this map. Nor does the scalar normalization
$-4096$ define an orientation or a vacuum. Primitive generators and
relations, chiral collision maps, protected state preparation and any
holographic reconstruction remain additional mathematical structures.
